\documentclass[11pt]{article}

\usepackage[margin=1in]{geometry}
\usepackage{amsmath,amssymb,amsfonts,bm}
\usepackage{graphicx}
\usepackage{xcolor}
\usepackage{float}
\usepackage{listings}
\usepackage{enumitem}
\usepackage{booktabs}
\usepackage{hyperref}
\usepackage{xurl}
\usepackage{fancyhdr}

% ---------- URL and Hyperlink setup ----------
\Urlmuskip=0mu plus 1mu
\hypersetup{
	breaklinks=true,
	colorlinks=true,
	linkcolor=blue,
	urlcolor=blue
}

% ---------- MATLAB code style ----------
\lstdefinestyle{matlabstyle}{
	language=Matlab,
	basicstyle=\ttfamily\small,
	numbers=left,
	numberstyle=\tiny,
	stepnumber=1,
	numbersep=8pt,
	frame=single,
	columns=fullflexible,
	keepspaces=true,
	showstringspaces=false,
	keywordstyle=\bfseries,
	commentstyle=\itshape\color{gray},
}

% ---------- Header setup ----------
\pagestyle{fancy}
\fancyhf{}
\fancyhead[L]{November 2, 2025}
\fancyhead[C]{ECE 517: Homework 2.1}
\fancyhead[R]{Scott Nguyen}
\renewcommand{\headrulewidth}{0.4pt}
\fancyfoot[C]{\thepage}

\begin{document}
	
	\noindent\textbf{Problem 1. Regularized MMSE Criterion}
	
	\noindent
	A variation of the MMSE criterion minimizes the norm of the weight vector $\mathbf{w}$. 
	This is a way to control the complexity of the structure. The corresponding function is
	
	\[
	\mathcal{L}(\mathbf{x}, \mathbf{w}) = \mathbb{E}[e^2] + \lambda \lVert \mathbf{w} \rVert^2
	\]
	
	\begin{enumerate}[label=\textbf{(\arabic*)}, leftmargin=1.2cm]
		\item Make the derivation of the closed solution for $\mathbf{w}$.
		\item Work out an iterative solution using the same technique as used in the Least Mean Squares algorithm.
		\item Comment and compare both solutions in a short conclusion section.
	\end{enumerate}
	
	\vspace{1em}
	\noindent\textbf{Solution.}
	
	We are given the regularized MMSE cost
	\[
	J(\mathbf{w}) = \mathbb{E}\big[(d - \mathbf{w}^\top \mathbf{x})^2\big] + \lambda \lVert \mathbf{w} \rVert^2,
	\]
	where $\mathbf{x}$ is the input vector, $d$ is the desired response, $\mathbf{w}$ is the weight vector we want to design, and $\lambda \ge 0$ is the regularization parameter that keeps the weights from growing too large.
	
	\medskip
	\noindent\textbf{(1) Closed-form solution}
	
	Start from the MSE part:
	\[
	\mathbb{E}\big[(d - \mathbf{w}^\top \mathbf{x})^2\big]
	= \mathbb{E}[d^2] - 2 \mathbf{w}^\top \mathbb{E}[d \mathbf{x}] + \mathbf{w}^\top \mathbb{E}[\mathbf{x}\mathbf{x}^\top]\mathbf{w}.
	\]
	To simplify notation, let
	\[
	\mathbf{p} = \mathbb{E}[d \mathbf{x}] \quad \text{and} \quad \mathbf{R} = \mathbb{E}[\mathbf{x}\mathbf{x}^\top].
	\]
	Then the cost becomes
	\[
	J(\mathbf{w}) = \mathbb{E}[d^2]
	- 2 \mathbf{p}^\top \mathbf{w}
	+ \mathbf{w}^\top \mathbf{R} \mathbf{w}
	+ \lambda \mathbf{w}^\top \mathbf{w}.
	\]
	Combine the quadratic terms:
	\[
	J(\mathbf{w}) = \mathbb{E}[d^2] - 2 \mathbf{p}^\top \mathbf{w}
	+ \mathbf{w}^\top (\mathbf{R} + \lambda \mathbf{I}) \mathbf{w}.
	\]
	
	Now differentiate w.r.t. $\mathbf{w}$:
	\[
	\nabla_{\mathbf{w}} J(\mathbf{w}) = -2 \mathbf{p} + 2(\mathbf{R} + \lambda \mathbf{I}) \mathbf{w}.
	\]
	Set this to zero to find the minimizer:
	\[
	(\mathbf{R} + \lambda \mathbf{I}) \mathbf{w} = \mathbf{p}.
	\]
	Assuming $\mathbf{R} + \lambda \mathbf{I}$ is invertible (it is if $\lambda>0$), we get
	\[
	\boxed{\mathbf{w}^\star = (\mathbf{R} + \lambda \mathbf{I})^{-1} \mathbf{p}.}
	\]
	
	\medskip
	\noindent\textbf{(2) Iterative solution (LMS-style)}
	
	For an online / adaptive version, we replace the expectations with the current sample $(\mathbf{x}(n), d(n))$. Define the instantaneous error
	\[
	e(n) = d(n) - \mathbf{w}^\top(n)\mathbf{x}(n).
	\]
	The instantaneous cost is
	\[
	J_n(\mathbf{w}) = e^2(n) + \lambda \lVert \mathbf{w} \rVert^2.
	\]
	Compute its gradient:
	\[
	\nabla_{\mathbf{w}} J_n(\mathbf{w}) = -2 e(n)\mathbf{x}(n) + 2 \lambda \mathbf{w}(n).
	\]
	A standard stochastic-gradient update with step size $\mu$ is
	\begin{align*}
		\mathbf{w}(n+1)
		&= \mathbf{w}(n) - \frac{\mu}{2} \nabla_{\mathbf{w}} J_n(\mathbf{w}(n)) \\
		&= \mathbf{w}(n) + \mu e(n)\mathbf{x}(n) - \mu \lambda \mathbf{w}(n).
	\end{align*}
	So the final iterative rule is
	\[
	\boxed{\mathbf{w}(n+1) = (1 - \mu \lambda)\,\mathbf{w}(n) + \mu e(n)\mathbf{x}(n).}
	\]
	This is exactly LMS but with a “weight decay” factor $(1 - \mu \lambda)$.
	
	\medskip
	\noindent\textbf{(3) Short comparison}
	
	\begin{itemize}
		\item \emph{Batch / closed form:} $\mathbf{w}^\star = (\mathbf{R} + \lambda \mathbf{I})^{-1} \mathbf{p}$ is precise but needs $\mathbf{R}$ and $\mathbf{p}$ and a matrix inverse.
		\item \emph{Iterative / LMS-like:} $\mathbf{w}(n+1) = (1 - \mu \lambda)\mathbf{w}(n) + \mu e(n)\mathbf{x}(n)$ works sample by sample and is good when data arrive over time.
		\item \emph{Effect of $\lambda$:} In both cases $\lambda$ penalizes large weights, makes the problem better conditioned, and usually improves generalization.
	\end{itemize}
	
\end{document}
